\abstract
This thesis presents the design, implementation, and deployment of RAM-USB,
a remotely accessible multi-user backup service, designed according to the zero-trust and
defense-in-depth principles.

The adopted methodology combines the spiral model for the definition and evolution of the requirements, 
the Test-Driven Development approach for the implementation, and the V-model for their verification. 
It also guarantees unidirectional traceability from requirement to code and bidirectional traceability from requirement
to commit and from requirement to test.

The work focuses in particular on the design of the authentication and dynamic access control
services, as well as on the solutions adopted to protect user credentials, critically analyzing
their advantages, disadvantages, and security implications. The system that implements them is
an n-tier client-server microservice architecture, organized into components dedicated to access,
registration, backup management, and metrics observability, deployed in Docker containers
which are in turn isolated in Kernel-based Virtual Machines and Linux Containers on one or more Proxmox hosts.

Internal communications take place via REST APIs for registration and login operations, and via 
MQTT for the publication of operational metrics; all of them remain confined to the mesh network coordinated by 
Headscale, mutually authenticated via mTLS with X.509 certificates issued by an internal 
Certificate-Authority.

In accordance with the zero-trust principle, every component revalidates the format and content
of every request, without assuming that the previous component has already done so. Following
authentication and authorization, the system dynamically modifies the access rules of the mesh
network, allowing the user to reach their own backup space via SFTP, isolated from that
of the other users by means of chroot. Backup files are deduplicated, compressed, and encrypted by
Restic before leaving the user's device, and their confidentiality therefore remains guaranteed by
client-side encryption.

The thesis concludes with a critical assessment of the security, confidentiality, and functionality properties
of the system in the scenario of full compromise of a single service. The analysis makes it possible to assess
the effectiveness of the isolation measures and of the defense-in-depth strategies adopted, while highlighting
their limitations and the properties that would be preserved even in the presence of a severe compromise
of the service.
